% chapters/52_coding_agents_ides.tex
% Part of: AI / LLM / Agents / Hardware Notes

\section{AI Coding Agents \& IDEs}

\subsection{Landscape overview}

% TODO

\subsection{Cursor: architecture \& edit model}

% TODO

\subsection{GitHub Copilot \& Codex}

% TODO

\subsection{Windsurf (Codeium)}

% TODO

\subsection{Antigravity}

% TODO

\subsection{Devin \& SWE-agent}

% TODO

\subsection{Aider \& CLI coding agents}

CLI-based coding agents run entirely in a terminal, wrapping an LLM with a persistent
conversation loop, file I/O tools, and shell execution.

\textbf{Aider} is an open-source CLI agent that maintains a \emph{repo map} (built with
ctags/tree-sitter) to give the model structural context without flooding the context window
with raw source. It applies model-generated edits via SEARCH/REPLACE blocks or unified-diff
format, and supports multi-model backends (GPT-4o, Claude, Gemini, etc.).

\textbf{Claude Code} (Anthropic, 2024--) is Anthropic's official CLI agent, built on the
Claude Agent SDK. Observations from direct inspection of live session JSONL logs:

\begin{itemize}
  \item \textbf{Session persistence.} Every conversation is stored as a newline-delimited
    JSON file at \texttt{\textasciitilde/.claude/projects/\{project-slug\}/\{session-uuid\}.jsonl}.
    The project slug is the working-directory path with \texttt{/} replaced by \texttt{-},
    so each repository gets its own project namespace.

  \item \textbf{Message graph.} Each JSONL entry carries \texttt{uuid} and \texttt{parentUuid},
    forming a linked list (or tree when the user branches). A \texttt{promptId} groups every
    event (user message, assistant chunks, tool results) that belongs to the same conversational
    turn. The boolean \texttt{isSidechain} marks branch points.

  \item \textbf{Event types.} Three kinds of records appear in the log:
    \begin{itemize}
      \item \texttt{queue-operation} (enqueue / dequeue) --- internal concurrency plumbing
        so messages do not collide when dispatched rapidly.
      \item \texttt{user} --- human messages \emph{and} tool-call results (see below).
      \item \texttt{assistant} --- model output, written incrementally as chunks arrive;
        a single logical response often spans several consecutive lines.
    \end{itemize}

  \item \textbf{Tool results are injected as \texttt{role: user}.} This is a direct
    consequence of how the Anthropic Messages API works: tool results are sent back to
    the model in a \texttt{user} turn, not as assistant output. Each result entry carries
    a \texttt{sourceToolAssistantUUID} linking it to the assistant message that issued
    the call, plus a \texttt{toolUseResult} object with \texttt{stdout}, \texttt{stderr},
    and an \texttt{interrupted} flag.

  \item \textbf{Extended thinking blocks.} When extended thinking is enabled, the
    assistant content array contains a \texttt{\{type: ``thinking''\}} block alongside
    the visible \texttt{\{type: ``text''\}} block. Each thinking block carries a
    cryptographic \texttt{signature} --- Anthropic's mechanism to verify that thinking
    content has not been tampered with between turns in a long tool-use loop.

  \item \textbf{Prompt-cache telemetry.} Each assistant message records
    \texttt{cache\_creation\_input\_tokens} (tokens written to the cache this turn) and
    \texttt{cache\_read\_input\_tokens} (tokens served from cache), broken out by
    TTL tier (\texttt{ephemeral\_5m} vs.\ \texttt{ephemeral\_1h}). This makes the cache
    hit rate observable as context grows.

  \item \textbf{Session metadata.} Every record stores \texttt{cwd}, \texttt{gitBranch},
    \texttt{version} (the Claude Code release string), and \texttt{userType}
    (\texttt{external} for human users). A human-readable \texttt{slug} field
    (e.g.\ \texttt{gleaming-spinning-nygaard}) appears from the first tool-call onward,
    matching the session's VM directory name. These fields enable session resumption and
    provide an audit trail.

  \item \textbf{Permission model.} Messages carry a \texttt{permissionMode} field
    (\texttt{default}, \texttt{full}, \texttt{restricted}) that gates which tools the
    agent may invoke autonomously, aligning with Anthropic's principle of minimal
    footprint for agentic actions.
\end{itemize}

\subsection{Claude Code remote control}

Claude Code's \textbf{remote control} feature allows an active terminal session to be
accessed and continued from a web browser or mobile device, without moving any code or
state to the cloud. Everything keeps running on the original machine; only the
conversation stream is proxied to the remote client.

\textbf{Workflow.}
\begin{enumerate}
  \item In the project directory, run \texttt{claude remote-control}. This starts a
    new session and prints a URL. Pressing \texttt{Space} generates a QR code that
    can be scanned to open the session directly in the Claude mobile app or browser.
  \item To hand off a session already in progress, type \texttt{/remote-control}
    inside the running session to get the same URL.
  \item On the remote device, tap the \emph{environment} button and select the
    machine by name. Multiple machines / sessions can be listed simultaneously.
\end{enumerate}

\textbf{Architecture properties.}
\begin{itemize}
  \item \textbf{Nothing moves to the cloud.} The agent's tools, MCP server
    connections, file system access, and shell environment remain on the originating
    machine. The remote device is a thin client that streams the conversation.
  \item \textbf{Session continuity.} Typing on the mobile app updates the terminal
    in real time and vice versa --- the shared state is the JSONL session log on the
    local machine.
  \item \textbf{Spawn mode.} From a remote device you can also create entirely new
    Claude Code sessions on the connected machine, enabling you to launch parallel
    agents without being physically present.
  \item \textbf{Authentication.} The URL is tied to the Claude account; no
    separate credentials are required beyond being logged into the same account on
    the remote device.
\end{itemize}

\textbf{Persistent remote control.} By default, remote control must be explicitly
enabled per session. In \texttt{/config} (Claude Code's configuration), enabling
\texttt{remote-control} for all sessions removes this requirement.

\begin{notebox}
  Remote control is available on Mac, Team, and Enterprise plans (rolling out to
  Pro/Max). It is a natural extension of the JSONL session log architecture: because
  every turn is already serialised to disk, surfacing that stream to a second client
  requires no additional server-side state.
\end{notebox}

\subsection{How AI IDEs use context (repo map, RAG)}

% TODO

\subsection{Diff application \& conflict resolution}

Coding agents must apply model-generated edits back to the repository. Two dominant
formats exist and the choice meaningfully affects benchmark performance:

\textbf{Whole-file (whole) format:}
The model rewrites the entire file with the change applied. Simple to parse and apply;
but wastes tokens and compute rewriting unchanged content. Produces more reliable edits
for small files; becomes impractical for large codebases.

\textbf{Diff (unified-diff / SEARCH-REPLACE) format:}
The model emits only the changed lines, either as a unified diff or as a
SEARCH/REPLACE block. Far more token-efficient; the apply step can fail if the
surrounding context the model assumed does not match the current file (merge
conflicts).

\textbf{Observed impact:} On the Aider Polyglot benchmark, the same base model
scored 8\% in diff format vs 16\% in whole format because the model had not been
fine-tuned on diff-style output --- it would hallucinate diff syntax. After
instruction-tuning on diff examples, performance converged and ultimately exceeded
the whole-format ceiling.

\begin{intuitionbox}[title=The analogy]
  Whole format = re-drawing the entire picture to add a cloud.
  Diff format = drawing just the cloud. Diff is smarter; it requires the model to
  reason precisely about which lines to change and what the surrounding context is.
\end{intuitionbox}

\textbf{Conflict resolution} arises when the model's assumed file state diverges from
the actual file (concurrent edits, stale context). Agents handle this by: (1) re-reading
the file before applying, (2) fuzzy matching the search block with edit distance
tolerance, or (3) falling back to whole-file replacement.

\subsection{Evaluation: SWE-bench, Aider Polyglot \& coding benchmarks}

\textbf{SWE-bench / SWE-bench Verified:}
\begin{itemize}
  \item Tasks are real GitHub issues from popular Python repositories; the agent must
    produce a patch that passes the associated unit tests.
  \item SWE-bench Verified filters to issues where human raters confirmed the test suite
    correctly validates the fix, removing noisy items.
  \item State-of-the-art agents (Devin, SWE-agent) score in the 40--50\% range as of
    mid-2024; fully autonomous resolution of hard issues remains challenging.
  \item Measures end-to-end capability: repo understanding, fault localisation,
    patch generation, and test passage --- all in one score.
\end{itemize}

\textbf{Aider Polyglot:}
\begin{itemize}
  \item Coding exercises in six languages (Python, JavaScript, TypeScript, Go, Rust,
    C++); evaluated in diff format and whole format separately.
  \item Respected as a real-world, language-diverse benchmark; frontier models like
    GPT-4-era scored around 18\% in diff format at launch.
  \item Benchmark non-determinism is significant: consecutive runs of the same
    checkpoint can vary by 2--3 percentage points (see Chapter~45 on statistical
    significance).
\end{itemize}

\textbf{Synthetic data for coding fine-tuning:}
Two prominent methods for generating instruction-tuning data for coding models:
\begin{itemize}
  \item \textbf{OSS-Instruct} (from MagicCoder): given a code snippet, prompt a
    stronger model to generate an instruction-response pair \emph{inspired by} the
    snippet but not identical to it. Produces diverse, naturally grounded examples.
  \item \textbf{Evol-Instruct:} iteratively evolve a seed instruction into more
    complex or diverse variants (``make it harder'', ``add an edge case'', ``translate
    to a different language''). Each evolved version becomes a new training example.
    The cloning/dancing-guy analogy: each iteration produces a slightly different
    version; after many rounds, the dataset is far richer than the seed.
\end{itemize}

\textbf{The training-data pipeline challenge} (observed empirically):
\begin{itemize}
  \item Raw code from GitHub commit messages varies wildly in quality; poor commit
    messages often indicate low-quality, uncommitted code artifacts.
  \item Validation harnesses (test runners that score model outputs) can \emph{pass
    garbage} if misconfigured --- a stricter harness catching more failures is
    progress, not regression.
  \item \textbf{Garbage in, garbage out:} training on lower-quality data can make a
    model worse than its base. Cleaning (whitespace normalisation, deduplication,
    license filtering, quality scoring) is as important as volume.
  \item Data sources used in practice: The Stack (60TB permissively-licensed code),
    GitHub (MIT/Apache-licensed repositories), synthetic generation from frontier
    models (with contamination checks).
\end{itemize}

